\section{The automaton and the easy direction}\label{sec:q10}

\begin{proposition}[completeness]\label{prop:complete}
A term with block $\bot$ has no head normal form.
\end{proposition}
\begin{proof}
If $t\twoheadrightarrow u$ with $u$ a head normal form, then $b(u)\ne\bot$ by
Lemma~\ref{lem:table}(ii), and $b(t)=b(u)$ by Lemma~\ref{lem:inv}.
\end{proof}

\citet[Fig.~5.14]{dots} give a bottom-up tree automaton with states
$Q=\{0,\dots,34\}$, initial state $0$ for $S$, transition table $\delta$, and one
rejecting state $\circledast=34$; it \emph{accepts} $t$ when the state $d(t)$, defined by
$d(S)=0$ and $d(u\,v)=\delta(d(u),d(v))$, is not $\circledast$. Question~10 asks whether
it accepts exactly the terms with a head normal form. Let $\varphi\colon B\to Q$ be the map
of \lean{HnfS.blockToDance}. Kernel evaluation checks that
\begin{equation}\label{eq:hom}
\varphi(\beta(a,a'))=\delta(\varphi(a),\varphi(a')),\qquad
\varphi(0)=0,\qquad
\varphi(a)=\circledast\iff a=\bot,
\end{equation}
and that $\varphi$ is onto. So the automaton is a homomorphic image of the block model,
and by induction on $t$, $d(t)=\varphi(b(t))$.

\begin{corollary}\label{cor:q10}
For $t\in\mathcal S$, the automaton rejects $t$ iff $b(t)=\bot$. Hence every rejected
term has no head normal form, and the automaton accepts exactly the terms with a head
normal form iff every h-fold term has one.
\end{corollary}
\begin{proof}
The first claim is $d(t)=\varphi(b(t))$ with~\eqref{eq:hom}. The second is
Proposition~\ref{prop:complete}. For the third, by the first claim acceptance is the
condition $b(t)\ne\bot$; one inclusion is Proposition~\ref{prop:complete}, and the other
is the statement that h-fold terms have head normal forms.
\end{proof}

The state $d(t)$ is therefore itself a reduction invariant. What remains is the direction
\citet{dots} call the challenge, soundness of the fold.
